% Options for packages loaded elsewhere
\PassOptionsToPackage{unicode}{hyperref}
\PassOptionsToPackage{hyphens}{url}
\documentclass[
  12pt,
  a4paper,
]{article}
\usepackage{xcolor}
\usepackage[margin=1in]{geometry}
\usepackage{amsmath,amssymb}
\setcounter{secnumdepth}{-\maxdimen} % remove section numbering
\usepackage{iftex}
\ifPDFTeX
  \usepackage[T1]{fontenc}
  \usepackage[utf8]{inputenc}
  \usepackage{textcomp} % provide euro and other symbols
\else % if luatex or xetex
  \usepackage{unicode-math} % this also loads fontspec
  \defaultfontfeatures{Scale=MatchLowercase}
  \defaultfontfeatures[\rmfamily]{Ligatures=TeX,Scale=1}
\fi
\usepackage{lmodern}
\ifPDFTeX\else
  % xetex/luatex font selection
\fi
% Use upquote if available, for straight quotes in verbatim environments
\IfFileExists{upquote.sty}{\usepackage{upquote}}{}
\IfFileExists{microtype.sty}{% use microtype if available
  \usepackage[]{microtype}
  \UseMicrotypeSet[protrusion]{basicmath} % disable protrusion for tt fonts
}{}
\makeatletter
\@ifundefined{KOMAClassName}{% if non-KOMA class
  \IfFileExists{parskip.sty}{%
    \usepackage{parskip}
  }{% else
    \setlength{\parindent}{0pt}
    \setlength{\parskip}{6pt plus 2pt minus 1pt}}
}{% if KOMA class
  \KOMAoptions{parskip=half}}
\makeatother
% definitions for citeproc citations
\NewDocumentCommand\citeproctext{}{}
\NewDocumentCommand\citeproc{mm}{%
  \begingroup\def\citeproctext{#2}\cite{#1}\endgroup}
\makeatletter
 % allow citations to break across lines
 \let\@cite@ofmt\@firstofone
 % avoid brackets around text for \cite:
 \def\@biblabel#1{}
 \def\@cite#1#2{{#1\if@tempswa , #2\fi}}
\makeatother
\newlength{\cslhangindent}
\setlength{\cslhangindent}{1.5em}
\newlength{\csllabelwidth}
\setlength{\csllabelwidth}{3em}
\newenvironment{CSLReferences}[2] % #1 hanging-indent, #2 entry-spacing
 {\begin{list}{}{%
  \setlength{\itemindent}{0pt}
  \setlength{\leftmargin}{0pt}
  \setlength{\parsep}{0pt}
  % turn on hanging indent if param 1 is 1
  \ifodd #1
   \setlength{\leftmargin}{\cslhangindent}
   \setlength{\itemindent}{-1\cslhangindent}
  \fi
  % set entry spacing
  \setlength{\itemsep}{#2\baselineskip}}}
 {\end{list}}
\usepackage{calc}
\newcommand{\CSLBlock}[1]{\hfill\break\parbox[t]{\linewidth}{\strut\ignorespaces#1\strut}}
\newcommand{\CSLLeftMargin}[1]{\parbox[t]{\csllabelwidth}{\strut#1\strut}}
\newcommand{\CSLRightInline}[1]{\parbox[t]{\linewidth - \csllabelwidth}{\strut#1\strut}}
\newcommand{\CSLIndent}[1]{\hspace{\cslhangindent}#1}
\setlength{\emergencystretch}{3em} % prevent overfull lines
\providecommand{\tightlist}{%
  \setlength{\itemsep}{0pt}\setlength{\parskip}{0pt}}
\usepackage{bookmark}
\IfFileExists{xurl.sty}{\usepackage{xurl}}{} % add URL line breaks if available
\urlstyle{same}
\hypersetup{
  hidelinks,
  pdfcreator={LaTeX via pandoc}}

\author{}
\date{}

\begin{document}

\section{Digital Twins for Software Engineers -- A Seven-Stage Deep
Research
Briefing}\label{digital-twins-for-software-engineers-a-seven-stage-deep-research-briefing}

\begin{quote}
Multi-perspective, corpus-grounded research. Method adapted from
hadufer/claude-storm (MIT) / Stanford OVAL STORM (Shao et al., NAACL
2024). Depth: quick (compact, \textasciitilde1000-word briefing) -
Perspectives: 4 + basic - Date: 2026-08-06
\end{quote}

\textbf{Scope.} This briefing targets a software engineer who has heard
``digital twin'' applied to manufacturing lines, buildings, and vehicles
and wants to know what, concretely, distinguishes it from ``simulation''
or ``monitoring dashboard'' as a software artifact. It covers
architecture, the software engineering lifecycle (build, test, sync),
and standardization; it does not cover twin-specific numerical modeling
methods (finite-element, CFD) or domain physics, which the corpus treats
as a separate concern from the software layer.

\subsection{Stage 1 -- Perspectives
assembled}\label{stage-1-perspectives-assembled}

Four grounded lenses plus a basic-fact pass: \textbf{the Practitioner}
(what building a twin surfaces beyond the papers), \textbf{the Academic}
(what the literature actually claims and where it disagrees),
\textbf{the Skeptic} (limitations the corpus admits to), and \textbf{the
Historian} (what a digital twin builds on). An Adoption/Incentives lens
was folded into the Practitioner lens, since the corpus discusses them
together.

\subsection{Stage 2 -- What the literature
claims}\label{stage-2-what-the-literature-claims}

A digital twin is consistently defined as a \emph{live, bidirectionally
synchronized} virtual representation of a physical asset -- distinct
from a plain simulation, which runs independently of any specific
physical instance, and from a digital shadow, which only flows data one
way {[}1{]}. Architecturally, most sources converge on a layered stack:
a physical asset, a data/connectivity layer, one or more models, and a
service layer exposing the twin's outputs {[}2{]} {[}3{]}. Model-driven
engineering (MDE) is repeatedly proposed as the natural fit for building
this stack, since twins already require explicit, versioned models of
structure and behavior {[}4{]}.

\subsection{Stage 3 -- What practitioners
encounter}\label{stage-3-what-practitioners-encounter}

Reference implementations expose engineering concerns the definitional
papers gloss over. \texttt{makeTwin} argues a general-purpose digital
twin platform needs first-class support for model building, data
processing, and low-code configuration for non-specialist developers,
not just a modeling notation {[}5{]}. \texttt{GreenhouseDT} frames a
twin as a layered, self-adaptive \emph{system}, not a static model, with
feedback loops that themselves need engineering and testing {[}6{]}.
TwinOps extends DevOps practice to model-based CPS engineering, treating
the twin's models as build artifacts that need CI pipelines of their own
{[}7{]}. That testing problem is concrete: cyber-physical systems teams
report continuous-integration practices that must account for
hardware-in-the-loop and simulation fidelity, not just code {[}8{]}, and
reproducibility work on the PiCar-X platform found that even a small
robotics twin needs disciplined environment and data versioning to get
repeatable results at all {[}9{]}.

\subsection{Stage 4 -- Where the corpus
disagrees}\label{stage-4-where-the-corpus-disagrees}

The clearest tension is scope: manufacturing-oriented sources treat
``digital twin'' as inseparable from a specific physical asset and its
live data feed {[}1{]} {[}5{]}, while standardization efforts (ISO
23247) formalize a more general reference architecture meant to apply
across an entity's whole lifecycle, independent of any one
implementation {[}10{]}. A second tension: the MDE literature treats a
digital twin primarily as a modeling artifact {[}4{]}, while the systems
literature (GreenhouseDT, TwinOps) treats it primarily as a running,
tested, deployed piece of software infrastructure {[}6{]} {[}7{]}. Both
are internally consistent; they disagree on which engineering discipline
owns the twin. Universal agreement: every source treats live
bidirectional synchronization with the physical asset as the
non-negotiable feature that separates a twin from a plain simulation.

\subsection{Stage 5 -- Core concepts for a software
engineer}\label{stage-5-core-concepts-for-a-software-engineer}

Practically, three things carry over directly from conventional software
engineering. First, architecture: a layered design -- asset,
connectivity, model, service -- with documented patterns for each layer
{[}11{]}. Second, lifecycle management: a twin's models are living
artifacts that need version control, CI, and testing pipelines just as
code does, and CPS-specific CI practice (hardware/simulation-in-
the-loop) is where most of the unfamiliar tooling lives {[}8{]} {[}7{]}.
Third, standardization: ISO 23247 gives a reference vocabulary and
structure worth adopting even for a one-off project, since it is the
point of convergence the corpus's otherwise-divergent architectures are
moving toward {[}10{]}. What does \emph{not} carry over unchanged is
testing: a twin's correctness depends on synchronization fidelity with a
live physical system, which unit and integration tests alone cannot
establish {[}9{]}.

\subsection{Stage 6 -- Synthesis and actionable
insight}\label{stage-6-synthesis-and-actionable-insight}

\textbf{Executive summary:} For a software engineer, a digital twin is
best understood as a distributed system with a live physical peer -- the
interesting engineering problems are synchronization, model lifecycle
management, and CI/CD for models rather than the physics itself.
\textbf{Key finding, ranked highest-reliability:} the
twin/simulation/shadow distinction by directionality of data flow is the
most consistently supported claim in the corpus {[}1{]}. \textbf{Hidden
connection,} visible only across perspectives: the MDE camp's emphasis
on models-as- artifacts and the systems camp's emphasis on tested,
deployed infrastructure are two halves of the same DevOps-for-models
problem that TwinOps names explicitly {[}7{]}. \textbf{Actionable
insight:} start any twin project by choosing where synchronization
fidelity is actually load-bearing for correctness, then design the CI
pipeline around verifying that, before investing in modeling notation.
\textbf{Frontier question:} which parts of CPS continuous-integration
practice generalize beyond the interview-based cases studied so far
{[}8{]}?

\subsection{Stage 7 -- Self peer review}\label{stage-7-self-peer-review}

Confidence is high (8/10) on the architecture and synchronization
claims, which recur across independent sources; lower (5/10) on the
CI/testing guidance, which rests on a small number of case studies.
Weakest link: the claim that ISO 23247 is the convergence point is
supported by standards papers but not yet by adoption data showing
practitioners actually converging on it -- worth verifying against
implementation surveys. Bias check: manufacturing and CPS sources
dominate; twins for pure software systems or IT infrastructure are
comparatively thin in this corpus. Missing sixth perspective: a
security/reliability lens on the live data channel between physical
asset and twin. Overall grade: B+ -- solid grounding on core definitions
and architecture, incomplete on testing practice at scale.

\subsection{References}\label{references}

\protect\phantomsection\label{refs}
\begin{CSLReferences}{0}{0}
\bibitem[\citeproctext]{ref-liu_review_2021}
\CSLLeftMargin{{[}1{]} }%
\CSLRightInline{M. Liu, S. Fang, H. Dong, and C. Xu, {``Review of
digital twin about concepts, technologies, and industrial
applications,''} \emph{Journal of Manufacturing Systems}, vol. 58, pp.
346--361, Jan. 2021, doi:
\href{https://doi.org/10.1016/j.jmsy.2020.06.017}{10.1016/j.jmsy.2020.06.017}.}

\bibitem[\citeproctext]{ref-ferko_architecting_2022}
\CSLLeftMargin{{[}2{]} }%
\CSLRightInline{E. Ferko, A. Bucaioni, and M. Behnam, {``Architecting
{Digital} {Twins},''} \emph{IEEE Access}, vol. 10, pp. 50335--50350,
2022, doi:
\href{https://doi.org/10.1109/ACCESS.2022.3172964}{10.1109/ACCESS.2022.3172964}.}

\bibitem[\citeproctext]{ref-heaton_platform_nodate}
\CSLLeftMargin{{[}3{]} }%
\CSLRightInline{L. Heaton, {``Platform {Stack} {Architectural}
{Framework}: {An} {Introductory} {Guide}.''}}

\bibitem[\citeproctext]{ref-michael_model-driven_2025}
\CSLLeftMargin{{[}4{]} }%
\CSLRightInline{J. Michael \emph{et al.}, {``Model-{Driven}
{Engineering} for {Digital} {Twins}: {Opportunities} and
{Challenges},''} \emph{Systems Engineering}, vol. 28, no. 5, pp.
659--670, 2025, doi:
\href{https://doi.org/10.1002/sys.21815}{10.1002/sys.21815}.}

\bibitem[\citeproctext]{ref-tao_maketwin_2024}
\CSLLeftMargin{{[}5{]} }%
\CSLRightInline{F. Tao \emph{et al.}, {``{makeTwin}: {A} reference
architecture for digital twin software platform,''} \emph{Chinese
Journal of Aeronautics}, vol. 37, no. 1, pp. 1--18, Jan. 2024, doi:
\href{https://doi.org/10.1016/j.cja.2023.05.002}{10.1016/j.cja.2023.05.002}.}

\bibitem[\citeproctext]{ref-kamburjan_greenhousedt_2024}
\CSLLeftMargin{{[}6{]} }%
\CSLRightInline{E. Kamburjan \emph{et al.}, {``{GreenhouseDT}: {An}
{Exemplar} for {Digital} {Twins},''} in \emph{Proceedings of the 19th
{International} {Symposium} on {Software} {Engineering} for {Adaptive}
and {Self}-{Managing} {Systems}}, in {SEAMS} '24. New York, NY, USA:
Association for Computing Machinery, Jun. 2024, pp. 175--181. doi:
\href{https://doi.org/10.1145/3643915.3644108}{10.1145/3643915.3644108}.}

\bibitem[\citeproctext]{ref-hugues_twinops_2022}
\CSLLeftMargin{{[}7{]} }%
\CSLRightInline{J. Hugues, J. Yankel, J. Hudak, and A. Hristozov,
{``Twinops: {Digital} twins meets devops,''} \emph{CARNEGIE-MELLON UNIV
PITTSBURGH PA, Tech. Rep.}, 2022, Accessed: Jan. 08, 2025. {[}Online{]}.
Available: \url{https://apps.dtic.mil/sti/trecms/pdf/AD1168471.pdf}}

\bibitem[\citeproctext]{ref-zampetti_continuous_2023}
\CSLLeftMargin{{[}8{]} }%
\CSLRightInline{F. Zampetti, D. Tamburri, S. Panichella, A. Panichella,
G. Canfora, and M. Di Penta, {``Continuous {Integration} and {Delivery}
{Practices} for {Cyber}-{Physical} {Systems}: {An} {Interview}-{Based}
{Study},''} \emph{ACM Trans. Softw. Eng. Methodol.}, vol. 32, no. 3, pp.
73:1--73:44, Apr. 2023, doi:
\href{https://doi.org/10.1145/3571854}{10.1145/3571854}.}

\bibitem[\citeproctext]{ref-barbie_toward_2024}
\CSLLeftMargin{{[}9{]} }%
\CSLRightInline{A. Barbie and W. Hasselbring, {``Toward
{Reproducibility} of {Digital} {Twin} {Research}: {Exemplified} with the
{PiCar}-{X}.''} arXiv, Aug. 2024. doi:
\href{https://doi.org/10.48550/arXiv.2408.13866}{10.48550/arXiv.2408.13866}.}

\bibitem[\citeproctext]{ref-shao_analysis_2023}
\CSLLeftMargin{{[}10{]} }%
\CSLRightInline{G. Shao, S. Frechette, and V. Srinivasan, {``An
{Analysis} of the {New} {ISO} 23247 {Series} of {Standards} on {Digital}
{Twin} {Framework} for {Manufacturing},''} American Society of
Mechanical Engineers Digital Collection, Sep. 2023. doi:
\href{https://doi.org/10.1115/MSEC2023-101127}{10.1115/MSEC2023-101127}.}

\bibitem[\citeproctext]{ref-tekinerdogan_systems_2020}
\CSLLeftMargin{{[}11{]} }%
\CSLRightInline{B. Tekinerdogan and C. Verdouw, {``Systems
{Architecture} {Design} {Pattern} {Catalog} for {Developing} {Digital}
{Twins},''} \emph{Sensors}, vol. 20, no. 18, p. 5103, Jan. 2020, doi:
\href{https://doi.org/10.3390/s20185103}{10.3390/s20185103}.}

\end{CSLReferences}

\end{document}
